% !TEX root = ../proyect.tex

\chapter{Definición de objetivos}\label{cap:objetivos}

\section{Objetivo general}\label{sec:objetivo-general}

El objetivo general de este Trabajo de Fin de Grado es desarrollar un procedimiento computacional para analizar series temporales desde la perspectiva de la dinámica no lineal y la predicción a corto plazo. El trabajo no parte de la idea de demostrar caos en una serie económica concreta, sino de construir una secuencia de herramientas que permita caracterizar dependencia temporal, reconstruir espacios de estados, comparar modelos predictivos y distinguir entre indicios de estructura y conclusiones no justificadas.

Para comprobar el comportamiento del procedimiento se utilizan dos contextos. El primero es sintético: el mapa logístico, como sistema determinista, no lineal y caótico conocido. El segundo es empírico: la volatilidad realizada intradía de Bitcoin, construida a partir de velas de cinco minutos del par \texttt{BTCUSDT}. Esta segunda parte permite estudiar qué ocurre cuando las herramientas se aplican a una serie financiera real, donde la señal aparece mezclada con ruido, heterocedasticidad y cambios de régimen.

Además del estudio técnico, el proyecto incluye un MVP web. Su finalidad es trasladar una parte del procedimiento a una herramienta interactiva capaz de cargar datos, construir variables y mostrar predicciones de volatilidad. El prototipo no sustituye al análisis completo, sino que actúa como demostración funcional de los modelos y artefactos generados durante el trabajo.

\section{Objetivos específicos}\label{sec:objetivos-especificos}

Para alcanzar el objetivo general se plantean los siguientes objetivos específicos:

\begin{enumerate}
    \item Presentar los conceptos necesarios para interpretar el trabajo, incluyendo volatilidad realizada, dependencia temporal, no linealidad, reconstrucción por retardos, exponentes de Lyapunov, entropía y predicción local.

    \item Comprobar el procedimiento sobre un sistema sintético conocido, en particular el mapa logístico, para observar su comportamiento cuando la dinámica subyacente es determinista y no lineal.

    \item Aplicar el procedimiento a una serie financiera real, construyendo una base de datos de volatilidad realizada intradía de Bitcoin a partir de velas de cinco minutos de \texttt{BTCUSDT}.

    \item Definir una variable principal de análisis basada en volatilidad realizada, separando la información disponible en cada instante del objetivo futuro utilizado para la evaluación predictiva.

    \item Caracterizar la serie empírica mediante herramientas descriptivas, análisis de estacionariedad, autocorrelación, espectro y recurrencia.

    \item Estudiar la presencia de dependencia temporal e indicios de no linealidad mediante filtrado lineal, análisis de residuos, contrastes complementarios y comparación con controles.

    \item Reconstruir espacios de estados a partir de retardos temporales, seleccionando parámetros mediante información mutua, falsos vecinos y el método de Cao.

    \item Cuantificar de forma exploratoria la dinámica reconstruida mediante dimensión de correlación, estimaciones de divergencia local y entropía de permutación.

    \item Evaluar la utilidad predictiva de la reconstrucción mediante métodos locales basados en vecinos y compararlos con modelos de referencia.

    \item Incorporar un modelo HAR-logRV como referencia multiescala de volatilidad realizada y como candidato práctico para su integración en el MVP.

    \item Diseñar e implementar un MVP web que permita cargar datos, calcular variables, ejecutar modelos y visualizar predicciones de volatilidad.

    \item Documentar las limitaciones metodológicas y prácticas del trabajo, especialmente en relación con la interpretación de la no linealidad, la predicción financiera y el alcance del prototipo.
\end{enumerate}

\section{Alcance del trabajo}\label{sec:alcance}

El alcance del trabajo se organiza en torno a una metodología y dos contextos de aplicación. El primero es sintético y utiliza el mapa logístico como sistema de control. El segundo es empírico y se centra en la volatilidad realizada de Bitcoin a horizonte horario. En este segundo caso, la variable estudiada no es el precio del activo ni la dirección futura del mercado, sino una medida de variabilidad construida a partir de retornos observados.

Aunque el procedimiento podría aplicarse a otras series temporales, en esta memoria solo se desarrolla completamente sobre el mapa logístico y sobre Bitcoin. Otros activos, otras frecuencias, otros horizontes de predicción o nuevas funciones sintéticas quedan fuera del alcance principal y se consideran posibles extensiones futuras.

El trabajo incluye una parte experimental y una parte software. La parte experimental comprende la construcción de la serie, su caracterización, la reconstrucción del espacio de estados, la cuantificación dinámica, la comparación con controles y la evaluación predictiva. La parte software consiste en un MVP web que permite utilizar una selección de modelos y artefactos obtenidos en el estudio técnico.

No se busca construir el modelo predictivo más complejo posible ni optimizar una estrategia de mercado. El interés está en comparar herramientas con distinto grado de complejidad y en analizar qué puede concluirse razonablemente al aplicar técnicas de dinámica no lineal a una serie financiera real.

\section{Limitaciones del trabajo}\label{sec:limitaciones-objetivos}

El trabajo presenta varias limitaciones que condicionan la interpretación de sus resultados. La primera procede del propio objeto de estudio: las series financieras reales son ruidosas, presentan cambios de régimen y no son estacionarias en sentido estricto. Esto dificulta separar con claridad dependencia lineal, no linealidad, heterocedasticidad y efectos externos no observados.

La segunda limitación afecta a las herramientas de dinámica no lineal. La reconstrucción del espacio de estados, la dimensión de correlación, las estimaciones de Lyapunov o la entropía de permutación pueden aportar información útil, pero no constituyen por sí solas una demostración de caos determinista en una serie económica real. En el mapa logístico se conoce la regla generadora; en Bitcoin solo pueden defenderse indicios y comparaciones empíricas.

La tercera limitación se refiere a la predicción. El trabajo evalúa volatilidad realizada futura, no rentabilidad ni dirección del precio. Por tanto, una mejora en error predictivo no implica automáticamente utilidad operativa en mercado ni debe interpretarse como una señal de compra o venta.

La cuarta limitación está relacionada con el MVP web. El prototipo demuestra la integración técnica de una parte del estudio, pero no incorpora todos los elementos necesarios para una herramienta financiera profesional. Su función es mostrar de forma interactiva los modelos desarrollados y sus salidas principales.

Por último, los resultados empíricos se limitan a Bitcoin, a la frecuencia de cinco minutos y al horizonte horario considerado. Cualquier generalización a otros activos, periodos o frecuencias requeriría repetir el procedimiento completo.

\section{Aportación principal del TFG}\label{sec:aportacion-principal}

La aportación principal del TFG consiste en desarrollar una metodología reproducible para analizar series temporales mediante herramientas de dinámica no lineal y predicción, y en trasladar parte de esa metodología a una implementación software utilizable. El trabajo no se limita a aplicar modelos predictivos de forma aislada, sino que construye una secuencia que parte de la explicación conceptual, se comprueba en un sistema sintético conocido y se aplica después a datos financieros reales.

Desde el punto de vista metodológico, la aportación está en adaptar herramientas de sistemas dinámicos a un contexto de series observadas, manteniendo una interpretación prudente. El procedimiento permite estudiar dependencia temporal, estructura residual, reconstrucción por retardos, cuantificación dinámica y predicción local, pero evita convertir esos indicios en afirmaciones no justificadas sobre caos determinista.

Desde el punto de vista software, la aportación consiste en organizar el análisis como un sistema reproducible y en desarrollar un MVP web que consume artefactos generados por el estudio técnico. Esta parte permite cerrar el trabajo desde la perspectiva de la ingeniería del software, conectando metodología, resultados y uso interactivo.

En conjunto, el TFG aporta una aproximación integrada: explica y aplica herramientas de dinámica no lineal, las comprueba sobre una función caótica conocida, las lleva a una serie financiera real y construye una aplicación que permite utilizar parte de los resultados obtenidos. Su contribución está en mostrar qué puede analizarse, qué debe compararse y qué límites deben respetarse al aplicar estas técnicas a datos reales.